\section{Study III: knowledge transfer from CT to radiograph classifiers with Swin transformers}
\label{sec:xm}
% Fourth analysis plan (SAP_v4.md with DESIGN_v4.md), hash-frozen in freeze_v4.json on 8 November 2026 before the first v4 training run.
% Numbers in Section sec:xm_results come from tables/v4_*.tex (tools/xm_tables.py) and results/final_xm_stats.json; do not retype them.
Studies I and II used CT only to render training images. Study III asks whether CT can teach a radiograph classifier directly: a CT
classifier is distilled \citep{hinton2015distilling} into X-ray models on projections of the same volumes \citep{angaitkar2026anatomy}, and
one transformer family serves both modalities so that their heat maps can be compared in the same patient. Five hypotheses were fixed in a
fourth analysis plan, hash-frozen on 8 November 2026 before any model of this study was trained.

\subsection{Data}
\label{sec:xm_data}
CT-RATE \citep{hamamci2026ctrate} provides non-contrast chest CT volumes with report-derived findings, and every Qure.ai projection is built
from one CT-RATE volume (Section~\ref{sec:data}). Six findings are shared: lung opacity, pleural effusion, atelectasis, cardiomegaly,
consolidation and lung lesion. The CT subset holds one reconstruction per patient, every validation and test patient of the projection split
and \num{1600} training patients drawn with a fixed seed (\num{2156} volumes: \num{1600} training, 259 validation, 297 test). Volumes that
CT-RATE lists as not showing the chest or not scanned head first and supine were excluded, together with their projections, leaving
\num{20533} training, 396 validation and 415 test projections. Each volume was reduced in the cloud environment to 24 axial slices over the
lung extent with three Hounsfield windows as channels (lung, mediastinum, bone), cropped to the body at 256\,px. CT-RATE is CC BY-NC-SA~4.0;
our design was stricter and kept everything derived from it in the cloud accounts except metrics and $32\times32$ heat-map frames, but in
practice teacher outputs, predictions and model weights were also transferred to the authors' computer (Section~\ref{sec:limitations}).
None of these is redistributed, and no CT image is shown. NIH ChestX-ray14 and
VinBigData are used with the splits of Study II. Pleural effusion, atelectasis, cardiomegaly and consolidation map to the classes of the
same name and lung lesion to NIH Mass or Nodule (VinBigData Nodule/Mass); these five are the exactly mapped shared findings. Lung opacity maps
only approximately to NIH Infiltration and is reported but not tested.

\subsection{Models and training}
\label{sec:xm_models}
All models use SwinV2-T (\texttt{swinv2\_tiny\_window8\_256}, ImageNet-1k weights; \citealp{liu2022swinv2,wightman2019timm}) with the
class-query decoder of Study II, a learning rate of $3\times10^{-4}$ with the encoder at a quarter of it, model selection on the validation
split and five seeds per model (Fig.~\ref{fig:architecture}c). The \emph{CT teacher} is a 2.5D network: the encoder processes each of the 24
slices with the three windows as colour channels, each slice yields $8\times8$ tokens with a two-dimensional sine code and a learned
slice-depth embedding, and the decoder attends over the tokens of all slices (20 epochs, four volumes per step); a volumetric Swin network did
not fit a T4 GPU at a useful batch size. The \emph{projection model} P0 learns the six findings from the projections (15 epochs, batch 32);
P1 adds, on the \num{1600} training projections whose volume is in the CT subset, a distillation and an alignment term. For a paired
projection $k$ with student logits $z_k$, teacher logits $t_k$ and pooled decoder outputs $e_k$ (student) and $e^{\mathrm{T}}_k$ (teacher),
\begin{equation}
\mathcal{L}=\mathcal{L}_{\mathrm{BCE}}(z,y)+\lambda_{\mathrm{kd}}\,T^{2}\,\mathcal{L}_{\mathrm{BCE}}\!\big(z_k/T,\,\sigma(t_k/T)\big)
+\lambda_{\mathrm{al}}\big(1-\cos(We_k,\,e^{\mathrm{T}}_k)\big),
\label{eq:kd}
\end{equation}
with $y$ the six finding labels, $\mathcal{L}_{\mathrm{BCE}}$ the binary cross-entropy averaged over findings, $\sigma$ the logistic
function, $T=2$,
$\lambda_{\mathrm{kd}}=1$, $\lambda_{\mathrm{al}}=0.5$ and $W$ a learned linear map; half of every P1 batch is drawn from the paired rows,
and P1 seed $k$ distils from the teacher of slot $k$. The \emph{radiograph classifiers} F$_\mathrm{in}$, F$_\mathrm{p0}$ and F$_\mathrm{p1}$
start from ImageNet, P0 and P1 (encoder and decoder loaded, heads re-initialised) and are trained on the 14 NIH diseases for 10 epochs with
the radiograph as the only input; the five DenseNet-121 radiograph-only runs of Study II are the convolutional reference (trained for 20
epochs at the full encoder learning rate, the pre-specified rule for batch-normalised encoders). Teacher seed 2 did
not learn (validation macro-AUROC 0.46--0.54 in every epoch, constant logits). Before any P1 model of seed 2 was trained we fixed a rule:
train seeds 5 and 6 with the unchanged recipe and give slot 2 to seed 5 unless its logits were constant (per-class standard deviation below
0.01), otherwise to seed 6. Seed 5 learned, so the teachers of slots 0--4 are seeds 0, 1, 5, 3 and 4 (Supplementary Table~S1).

\subsection{Explanations, cross-modal agreement and hallucination tests}
\label{sec:xm_xai}
For the radiograph classifiers we computed Grad-CAM++ on the last encoder stage \citep{chattopadhay2018gradcampp}, Integrated Gradients
\citep{sundararajan2017axiomatic} and the attention of the last decoder layer, scored with the pointing game on the 984 NIH expert boxes
\citep{zhang2018topdown} at the 15\,px tolerance of Study II. For each CT test scan positive for a finding, the teacher's slice-wise
Grad-CAM++ maps were stacked and summed along the anteroposterior axis, giving a frontal map in the geometry of the paired projection. CT and
X-ray maps were cropped to their own lung boxes, resized to $32\times32$, scaled to sum to one and averaged over the five seed pairs.
Agreement is the Pearson correlation of the matched maps, and the null correlates the CT map with the X-ray maps of the other scans positive
for the same finding. A pre-specified orientation check on the attenuation frames alone (correlation 0.472 as they were, 0.435 mirrored)
left the maps unmirrored.

We call a report of a finding that the image does not contain a hallucination. On the \num{9861} NIH test images without any finding, with
seed-mean probabilities and per-class thresholds at 90\,\% specificity on the validation split, the hallucination rate is the share of
images in which some class moves from at most to above its threshold after a perturbation that adds no disease: rib emphasis (the radiograph
plus half of the bone image of the Study~I suppressor), Gaussian noise ($\sigma=0.02$) or a gamma of 0.8. Counterfactual nodules from the
audit library (Section~\ref{sec:audit}) were pasted into lung fields of images whose Mass and Nodule scores were at most their thresholds;
sensitivity is the share in which either score crosses it. Off-target attribution is the share of Grad-CAM++ mass outside the lung fields for
positive predictions of lung findings.

\subsection{Hypotheses and statistics}
\label{sec:xm_stats}
A model's prediction is the mean of its seeds' probabilities. H11 (primary) compares F$_\mathrm{p1}$ with F$_\mathrm{in}$ in the mean NIH
test AUROC of the five shared findings; H12 compares F$_\mathrm{p1}$ with F$_\mathrm{p0}$ and F$_\mathrm{p0}$ with F$_\mathrm{in}$; H13
compares F$_\mathrm{in}$ with DenseNet-121 in the 14-disease macro-AUROC; that reference had been evaluated on the NIH test split in Study II
about ten hours before this plan was frozen, so its score was known when H13 was specified (no Swin model existed yet). These contrasts use a paired patient-cluster bootstrap
(\num{10000} resamples, two-sided percentile $p$-values, Eq.~\eqref{eq:bootp}); H11 and the two H12 contrasts form one Holm family. H14 is
the mean per-map difference between matched and mismatched correlation (scan-cluster bootstrap) and H15 the mean per-image difference in the
rib-emphasis flip between F$_\mathrm{p1}$ and F$_\mathrm{in}$ (patient-cluster bootstrap). The CT teacher, the projection models and their
zero-shot transfer, per-class DeLong tests \citep{delong1988auc}, conformal risk control of the false-negative rate at 0.10
\citep{angelopoulos2022crc} and the pointing game are descriptive; the grouping of the DeLong tests and the pointing-game summary were fixed
after the hypothesis results were known.
